\chapter{Die Beweise}
\label{ch:anhang-beweise}

Die Kapitel dieses Buchs erzählen die Beweisideen; hier stehen die
Beweise selbst, nach Kapiteln geordnet. Notation wie im Haupttext;
$E_n$ bezeichnet die Codierung $V_n^{-1}$, $\mu(x)$ den
Minimalhorizont, $\sigma(P)=\max_{c_d>0}(d+c_d)$ (mit
$\sigma(0)=0$) den Strukturhorizont.

\section*{Zu Kapitel 3: die Grenzraumsätze}

\begin{satz}[Satz~\ref{satz:grenzen}, vollständig]
$\varinjlim(H_n,Z_n)\cong\Struk$ und
$\varinjlim(H_n,L_n)\cong\Nn$.
\end{satz}

\begin{proof}
\emph{Strukturtreu.} Für $m\ge n$ sei $Z_{n,m}$ das Voranstellen von
$m-n$ Nullen; offenbar $Z_{m,k}\circ Z_{n,m}=Z_{n,k}$, das System
ist gerichtet. Im direkten Limes gelten $a\in H_n$ und $b\in H_m$
(o.B.d.A.\ $n\le m$) als gleich, wenn sie in einem gemeinsamen
späteren Horizont dasselbe Bild haben, also
$0^{k-n}a=0^{k-m}b$, was äquivalent zu $b=0^{m-n}a$ ist. Die
Klassen sind also Ziffernfolgen modulo führender Nullen. Die
Abbildung $[a]\mapsto P_a$ ist wohldefiniert (eine führende Null
fügt den Koeffizienten $0$ am höchsten Grad hinzu), injektiv (die
Ziffernfolge in $H_n$ ist die auf Länge $n$ aufgefüllte, von rechts
gelesene Koeffizientenfolge; gleiche Polynome erzwingen also
$b=0^{m-n}a$) und surjektiv (jedes $P\in\Struk$ ist ab
$H_{\sigma(P)}$ darstellbar).

\emph{Werttreu.} Sei $L_{n,m}=E_m\circ V_n$. Wegen
$V_m\circ E_m=\mathrm{id}$ gilt
$L_{m,k}\circ L_{n,m}=E_k\circ V_n=L_{n,k}$. Zwei Elemente werden
genau dann identifiziert, wenn $E_k(V_n(a))=E_k(V_m(b))$ für ein
$k$, und da $E_k$ injektiv ist, genau bei $V_n(a)=V_m(b)$: Die
Klassen sind die Fasern der Wertabbildung. $[a]\mapsto V_n(a)$ ist
wohldefiniert, injektiv und surjektiv (jedes $x$ tritt als
$V_{\mu(x)}(E_{\mu(x)}(x))$ auf).
\end{proof}

\section*{Zu Kapitel 3: die dritte Treppe}

\begin{satz}[Cantor-Koordinate und Rechts-Erweiterung]
Mit $F(a)=\sum_i a_i/(i+1)!$ gilt $V_n(a)=(n+1)!\cdot F(a)$. Das
Rechts-Anhängen $R(a_1,\ldots,a_n)=(a_1,\ldots,a_n,0)$ ist
bruchtreu, skaliert den Wert exakt mit $n+2$, und
$\varinjlim(H_n,R_n)$ ist die Menge der endlichen
Fakultätsbrüche in $[0,1)$.
\end{satz}

\begin{proof}
Termweise ist $a_i\,(n+1)!/(i+1)!=(n+1)!\cdot a_i/(i+1)!$, also die
Faktorisierung. Beim Rechts-Anhängen behalten alle Ziffern Position
und Nenner, die neue Null trägt nichts bei: $F$ invariant; der Wert
ist $(n+2)!\,F=(n+2)\cdot(n+1)!\,F$. Im Limes werden genau
Ziffernfolgen identifiziert, die sich um angehängte Nullen
unterscheiden; jede Klasse trägt genau einen endlichen Bruch, und
jeder endliche Fakultätsbruch tritt auf.
\end{proof}

\section*{Zu Kapitel 4: die Horizontableitung}

\begin{satz}[$\Delta$-Kalkül]
Die Wertänderung bei strukturtreuer Erweiterung ist
$\Delta P(n+1)$ mit $\Delta P(X)=P(X{+}1)-P(X)$ und
$\Delta\,\ff Xd=d\,\ff X{d-1}$; eine führende Null ist genau dann
wertneutral, wenn $P$ konstant ist.
\end{satz}

\begin{proof}
Unter $Z$ bleibt $P$ erhalten und der Auswertungspunkt wandert von
$n{+}1$ zu $n{+}2$; die Änderung ist per Definition
$\Delta P(n{+}1)$. Zur Differenzformel:
$\ff{(X{+}1)}d=(X{+}1)\,X\cdots(X{-}d{+}2)$ und
$\ff Xd=X\cdots(X{-}d{+}2)\,(X{-}d{+}1)$ haben den gemeinsamen
Faktor $X(X{-}1)\cdots(X{-}d{+}2)=\ff X{d-1}$; die verbleibenden
Faktoren sind $X{+}1$ bzw.\ $X{-}d{+}1$, ihre Differenz ist $d$,
also $\Delta\,\ff Xd=d\,\ff X{d-1}$.
Wertneutralität: $\Delta P=0$ für alle Auswertungspunkte im
zulässigen Bereich erzwingt, da Polynome mit unendlich vielen
Nullstellen verschwinden, $\Delta P\equiv0$, also $P$ konstant;
umgekehrt klar. Für den einzelnen Schritt: $\Delta P(n{+}1)=0$ mit
nichtnegativen Koeffizienten von $\Delta P$ (alle
Basisdifferenzen haben nichtnegative Koeffizienten) und
Auswertungspunkt im positiven Bereich -- wegen
$\deg\Delta P\le n-2$ sind alle $\ff{(n{+}1)}d$ dort strikt
positiv -- erzwingt ebenfalls $\Delta P=0$.
\end{proof}

\section*{Zu Kapitel 4: das Horizontkalkül}

Wir schreiben $P\preceq Q$, wenn jeder Koeffizient von $P$ in der
fallenden Fakultätsbasis höchstens der entsprechende Koeffizient von
$Q$ ist, und $M_r(X)=\sum_{d=0}^{r-1}(r{-}d)\,\ff Xd$ für die maximal
gefüllte Gestalt des Horizonts $r$.

\begin{satz}[Extremalprinzip; Register S39]
Sei $T$ ein $k$-stelliger Operator auf dem Strukturhalbring, der in
jedem Argument koeffizientenmonoton ist (aus $P\preceq P'$ folgt
$T(\ldots,P,\ldots)\preceq T(\ldots,P',\ldots)$). Dann gilt
\[
\Gamma_T(r_1,\ldots,r_k)
:=\max\{\sigma(T(P_1,\ldots,P_k)) \mid \sigma(P_i)\le r_i\}
=\sigma\bigl(T(M_{r_1},\ldots,M_{r_k})\bigr).
\]
Jeder aus Addition, Multiplikation und $\Delta$ zusammengesetzte
Operator ist koeffizientenmonoton.
\end{satz}

\begin{proof}
$\sigma(P)=\max_{c_d>0}(d+c_d)$ ist monoton bezüglich $\preceq$, und
$\sigma(P)\le r$ ist äquivalent zur koeffizientenweisen Schranke
$c_d\le r-d$, also zu $P\preceq M_r$. Aus der Monotonie von $T$
folgt $T(P_1,\ldots,P_k)\preceq T(M_{r_1},\ldots,M_{r_k})$ und damit
$\sigma(T(P_1,\ldots,P_k))\le\sigma(T(M_{r_1},\ldots,M_{r_k}))$;
wegen $\sigma(M_{r_i})=r_i$ wird die Schranke angenommen. Zur
Zusatzaussage: Addition und $\Delta$ sind koeffizientenweise linear
mit nichtnegativen Konstanten ($\Delta\,\ff Xd=d\,\ff X{d-1}$); für
die Multiplikation sind die Linearisierungskoeffizienten
$\binom pj\binom qj\,j!$ der Produktformel nichtnegativ, und die
Bilinearität überträgt die Monotonie. Kompositionen monotoner
Operatoren sind monoton.
\end{proof}

\begin{satz}[Differenzüberlauf; Register S40]
Für $r\ge2$ gilt
$\Gamma_\Delta(r)=\bigl\lfloor(r+1)^2/4\bigr\rfloor-1$,
angenommen bei der Ein-Ziffern-Gestalt $P=(r{-}d)\,\ff Xd$ mit
$d=\lfloor(r+1)/2\rfloor$.
\end{satz}

\begin{proof}
Nach dem Extremalprinzip ist $\Gamma_\Delta(r)=\sigma(\Delta M_r)$.
$\Delta M_r$ hat die Koeffizienten $b_k=(k{+}1)(r{-}k{-}1)$ für
$k=0,\ldots,r{-}2$, also
\[
\sigma(\Delta M_r)
=\max_{0\le k\le r-2}\bigl[k+(k{+}1)(r{-}k{-}1)\bigr]
=\max_{0\le k\le r-2}(k{+}1)(r{-}k)\;-\;1 .
\]
Mit $u=k{+}1\in\{1,\ldots,r{-}1\}$ ist $u(r{+}1{-}u)$ das Produkt
zweier Zahlen mit fester Summe $r{+}1$, maximal bei
$u=\lfloor(r{+}1)/2\rfloor$ mit Wert
$\lfloor(r{+}1)^2/4\rfloor$; der Bereich ist für $r\ge2$ nicht leer.
Die Ein-Ziffern-Gestalt $P=(r{-}d)\,\ff Xd$ hat $\sigma(P)=r$ und
$\sigma(\Delta P)=d-1+d(r{-}d)=d(r{+}1{-}d)-1$, maximal beim selben
$d$ mit demselben Wert. (Exhaustiv bestätigt für alle $(r{+}1)!$
Gestalten bis $r=8$: Werte $1,3,5,8,11,15,19$.)
\end{proof}

\begin{satz}[Kompositionsabschluss; Register S41]
Für $P,Q$ mit nichtnegativen ganzzahligen Koeffizienten in der
fallenden Fakultätsbasis hat auch $P\circ Q$ nichtnegative
ganzzahlige Koeffizienten. Die Komposition ist in beiden Argumenten
koeffizientenmonoton; nach dem Extremalprinzip gilt
$\Gamma_\circ(r,s)=\sigma(M_r\circ M_s)$.
\end{satz}

\begin{proof}
Wegen $P\circ Q=\sum_d c_d(P)\,(Q)^{\underline d}$ mit
$(Q)^{\underline d}=Q(Q{-}1)\cdots(Q{-}d{+}1)$ genügt es,
$(Q)^{\underline d}$ als zulässige Gestalt nachzuweisen. Schreibe
$Q=\sum_k a_k\ff Xk$ und wähle paarweise disjunkte Mengen $T_k$ mit
$|T_k|=a_k$ (\emph{Schablonen der Stelligkeit $k$}). Eine
\emph{$Q$-Struktur} auf $[x]$ sei ein Paar $(t,\varphi)$ aus einer
Schablone $t\in T_k$ und einer Injektion
$\varphi\colon[k]\hookrightarrow[x]$; ihre Anzahl ist
$\sum_k a_k\,\ff xk=Q(x)$. Also zählt $Q(x)^{\underline d}$ die
$d$-Tupel paarweise verschiedener $Q$-Strukturen auf $[x]$.

Klassifiziere die Tupel nach der Vereinigung $U\subseteq[x]$ der
Bilder aller beteiligten Injektionen. Die Anzahl $N_m$ der Tupel,
deren Bildvereinigung eine feste $m$-Menge ganz ausschöpft, hängt
nur von $m$ ab, und für alle $x\in\Nn$ gilt die Polynomidentität
$Q(x)^{\underline d}=\sum_m N_m\binom xm$. Die symmetrische Gruppe
$S_m$ wirkt auf den ausschöpfenden Tupeln durch Umbenennung der
Grundmenge, $\pi\cdot(t,\varphi)=(t,\pi\circ\varphi)$, und diese
Wirkung ist \emph{frei}: Fixiert $\pi$ jedes Tupelglied, so fixiert
$\pi$ jedes Bild punktweise, wegen der Ausschöpfung also ganz $[m]$,
d.\,h.\ $\pi=\mathrm{id}$. (Schablonen der Stelligkeit 0 stören
nicht: Sie tragen nichts zur Vereinigung bei.) Also $m!\mid N_m$,
und $(Q)^{\underline d}=\sum_m(N_m/m!)\,\ff Xm$ hat Koeffizienten in
$\Nn$ -- Positivität und Ganzzahligkeit in einem Zug.

Monotonie in $P$: $\Nn$-Linearkombination der
$(Q)^{\underline d}$. Monotonie in $Q$: Aus $a_k\le a_k'$ folgt
$T_k\subseteq T_k'$; jedes Tupel von $Q$-Strukturen ist eines von
$Q'$-Strukturen mit gleicher Bildvereinigung, also $N_m\le N_m'$
termweise. (Abschluss und Extremalprinzip rechnerisch exhaustiv
bestätigt für $r,s\le4$, der Abschluss zusätzlich für 300
Zufallspaare bis $\sigma=8$.)
\end{proof}

\begin{satz}[Shift-Überlauf; Register S42]
Für $E\,P(X)=P(X{+}1)$ gilt
$\Gamma_E(r)=r+\lfloor r^2/4\rfloor=\Gamma_\Delta(r{+}1)$.
\end{satz}

\begin{proof}
$E$ ist linear mit nichtnegativen Konstanten
($\ff{(X{+}1)}d=\ff Xd+d\,\ff X{d-1}$), also greift das
Extremalprinzip: $\Gamma_E(r)=\sigma(E\,M_r)$ mit Koeffizienten
$c_k=(r{-}k)+(k{+}1)(r{-}k{-}1)$ für $k\le r{-}2$ und $c_{r-1}=1$;
somit $k+c_k=r+(k{+}1)(r{-}k{-}1)$, und das Produkt zweier Zahlen
mit Summe $r$ ist maximal $\lfloor r^2/4\rfloor$. Die zweite
Gleichheit ist die Identität
$r+\lfloor r^2/4\rfloor=\lfloor(r{+}2)^2/4\rfloor-1$.
\end{proof}

\begin{satz}[Hori-\#P-Satz; Register S43]
Ist $f\in\#\mathrm P$ und $P$ eine Gestalt (nichtnegative
ganzzahlige Koeffizienten in der fallenden Fakultätsbasis), so ist
$P\circ f\in\#\mathrm P$.
\end{satz}

\begin{proof}
Sei $f(x)$ die Anzahl der Zeugen von $x$ unter einem polynomiellen
Verifizierer $V$, und $P=\sum_{d=0}^m c_d\ff Xd$. Neue Maschine auf
Eingabe $x$: Rate $d\in\{0,\ldots,m\}$, rate eine Variante
$\ell\in\{1,\ldots,c_d\}$, rate $d$ Zeugenkandidaten
$w_1,\ldots,w_d$; akzeptiere genau dann, wenn $V$ alle $w_i$
akzeptiert und die $w_i$ paarweise verschieden sind. Da $m$ und
alle $c_d$ feste Konstanten sind, bleibt die Maschine polynomiell,
und ihre Anzahl akzeptierender Pfade ist exakt
$\sum_d c_d\,f(x)^{\underline d}=P(f(x))$, denn
$f(x)^{\underline d}$ zählt die geordneten $d$-Tupel paarweise
verschiedener Zeugen. (Die Aussage folgt wegen
$\ff Xd=d!\binom Xd$ auch aus der bekannten Charakterisierung der
\emph{binomial-good}-Polynome als \#P-Abschlussoperationen,
Ikenmeyer--Pak, FOCS 2022; der obige Beweis ist selbständig.)
\end{proof}

\section*{Zu Kapitel 5: die beiden Halbringe und die Dualität}

\begin{satz}[Satz~\ref{satz:a3}, vollständig]
Mit $(n,x)\hplus(m,y)=(\max(n,m,\muh(x{+}y)),x{+}y)$ und
$(n,x)\htimes(m,y)=(\muh(xy),xy)$ ist $\Hori$ ein kommutativer
Halbring ohne Eins; $(0,0)$ ist neutral und absorbierend;
$h\htimes(1,1)=\KV(h)$; die wertkompakten Elemente bilden einen zu
$(\Nn,+,\cdot)$ isomorphen Unterhalbring mit Eins $(1,1)$.
\end{satz}

\begin{proof}
Vorbemerkung: $\muh$ ist monoton. \emph{Assoziativität von
$\hplus$:} Beide Klammerungen von $(n,x)$, $(m,y)$, $(k,z)$ liefern
den Wert $x{+}y{+}z$ und die Horizonte
$\max(n,m,k,\muh(x{+}y),\muh(x{+}y{+}z))$ bzw.\ mit $\muh(y{+}z)$;
wegen $x{+}y\le x{+}y{+}z$ und $y{+}z\le x{+}y{+}z$ absorbiert
$\muh(x{+}y{+}z)$ die inneren Terme. \emph{Assoziativität von
$\htimes$:} beide Seiten $(\muh(xyz),xyz)$.
\emph{Distributivität:} Links steht
$(n,x)\htimes\bigl((m,y)\hplus(k,z)\bigr)
=(\muh(xy{+}xz),\,xy{+}xz)$, rechts
$(\max(\muh(xy),\muh(xz),\muh(xy{+}xz)),\,xy{+}xz)$; wegen
$xy\le xy{+}xz$ und $xz\le xy{+}xz$ absorbiert $\muh(xy{+}xz)$
auch hier die inneren Terme.
\emph{Neutral/absorbierend:} $(n,x)\hplus(0,0)=(\max(n,0,\muh(x)),x)
=(n,x)$ wegen $n\ge\muh(x)$; $(n,x)\htimes(0,0)=(0,0)$.
\emph{Keine Eins:} Ein neutrales $(k,e)$ erforderte $xe=x$ für alle
$x$, also $e=1$, und $\muh(x)=n$ für alle Elemente -- das scheitert
an $(1,0)$. \emph{Projektor:}
$(n,x)\htimes(1,1)=(\muh(x),x)=\KV(n,x)$. \emph{Eck:} Für
wertkompakte Operanden ist
$\max(\muh(x),\muh(y),\muh(x{+}y))=\muh(x{+}y)$; unter $\htimes$ ist
das Ergebnis per Definition wertkompakt; $(1,1)$ ist dort neutral,
und $x\mapsto(\muh(x),x)$ ist der Isomorphismus.
\end{proof}

\begin{lemma}[$\sigma$-Additionsmonotonie]\label{lem:app-sigma}
$\sigma(P{+}Q)\ge\max(\sigma P,\sigma Q)$.
\end{lemma}

\begin{proof}
An der Maximalstelle $d^*$ von $P$ ist
$c_{d^*}(P{+}Q)\ge c_{d^*}(P)>0$, also
$\sigma(P{+}Q)\ge d^*+c_{d^*}(P{+}Q)\ge\sigma(P)$.
\end{proof}

\begin{satz}[Satz~\ref{satz:dualitaet}, vollständig]
Mit $(n,P)\boxplus(m,Q)=(\max(n,m,\sigmah(P{+}Q)),P{+}Q)$ und
$(n,P)\boxtimes(m,Q)=(\sigmah(PQ),PQ)$ ist $\Hori$ ein kommutativer
Halbring ohne Eins; $(0,0)$ ist neutral und absorbierend;
$h\boxtimes(1,1)=\KS(h)$; $\KS$ ist Homomorphismus beider
Operationen; die strukturkompakten Elemente bilden einen zu
$\Struk$ isomorphen Unterhalbring.
\end{satz}

\begin{proof}
Wörtlich spiegelbildlich zum vorigen Beweis, mit
Lemma~\ref{lem:app-sigma} anstelle der $\muh$-Monotonie. Zusätzlich
die Homomorphie von $\KS$: für $\boxplus$ ist
$\KS(a)\boxplus\KS(b)=(\max(\sigmah P,\sigmah Q,\sigmah(P{+}Q)),P{+}Q)
=(\sigmah(P{+}Q),P{+}Q)=\KS(a\boxplus b)$ nach dem Lemma; für
$\boxtimes$ sind beide Seiten $(\sigmah(PQ),PQ)$. Kein
$\boxtimes$-neutrales Element: es erforderte $PE=P$, also $E=1$,
und $\sigmah(P)=n$ für alle Elemente -- das scheitert an $(2,1)$,
denn $\sigmah(1)=1$.
\end{proof}

\begin{satz}[Produktprinzip, Satz~\ref{satz:produktprinzip}]
Jede kommutative Halbringstruktur $(\uplus,\odot)$ auf den
Exzessen liefert via
$(x,\varepsilon)\oplus(y,\delta)=(x{+}y,\varepsilon\uplus\delta)$,
$(x,\varepsilon)\otimes(y,\delta)=(xy,\varepsilon\odot\delta)$
eine gesetzestreue Arithmetik auf $\Hori$.
\end{satz}

\begin{proof}
Die Abbildung $(n,x)\mapsto(x,\,n-\muh(x))$ ist eine Bijektion
$\Hori\to\Nn\times\Nn$ (die Bedingung $n\ge\muh(x)$ ist genau
$\varepsilon\ge0$; die Umkehrung ist
$(x,\varepsilon)\mapsto(\muh(x){+}\varepsilon,x)$). Die Operationen
sind komponentenweise definiert, mit $(\Nn,+,\cdot)$ im Wertfaktor;
ein Produkt zweier kommutativer Halbringe erfüllt alle Gesetze
komponentenweise. Die Selektoren sind gültig, da der
Ergebnishorizont $\muh(\text{Wert})+\text{Exzess}\ge\muh$ ist.
\end{proof}

\begin{satz}[Kompaktifizierungsmonoid]
$\KS$ erhält Wertkompaktheit. Folglich gilt
$\KV\KS\KV=\KS\KV$, und das von $\KS,\KV$ erzeugte
Transformationsmonoid hat genau sechs Elemente -- $\mathrm{id}$,
$\KS$, $\KV$, $\KS\KV$, $\KV\KS$ und $\KS\KV\KS$ --, von denen
alle außer $\KV\KS$ idempotent sind.
\end{satz}

\begin{proof}
\emph{Stabilität:} Sei $(m,P)$ wertkompakt; das Element $(0,0)$
wird von $\KS$ fixiert, sei also $m\ge1$ und damit
$P(m{+}1)\ge m!$. Sei $s=\sigmah(P)$; angenommen, es gälte
$P(s{+}1)<s!$. Für $s\le t<m$ wächst
jeder Term $c_d\ff Xd$ von $t{+}1$ auf $t{+}2$ um den Faktor
$(t{+}2)/(t{+}2{-}d)$, wegen $d\le s{-}1\le t{-}1$ also um
höchstens $(t{+}2)/3$; damit
$P(t{+}2)<t!\,(t{+}2)/3\le(t{+}1)!$, und Induktion liefert
$P(m{+}1)<m!$ -- Widerspruch.
\emph{Relation:} Für jedes $h$ ist $\KV(h)$ wertkompakt, nach der
Stabilität auch $\KS(\KV(h))$, also fixiert $\KV$ dieses Element.
\emph{Sechs Elemente:} Jedes alternierende Wort der Länge $\ge4$
enthält $\KV\KS\KV$ und reduziert; Quadrate reduzieren per
Idempotenz -- höchstens sechs. Paarweise Verschiedenheit durch
explizite Trennzeugen (endliche Enumeration); insbesondere trennt
$h=(6,6)$: $\KV\KS(h)=(3,6)$, aber $\KS\KV\KS(h)=(2,5)$. Daraus
$(\KV\KS)^2=\KS\KV\KS\ne\KV\KS$: nicht idempotent; die Idempotenz
der übrigen fünf rechnet man mit der Relation nach.
\end{proof}

\section*{Zu Kapitel 5: das Tiefengesetz}

\begin{satz}[Satz~\ref{satz:tiefengesetz}, vollständig; Register S33]
Für jede Schale $m\ge3$ gilt
$\min\{\delta(h)\mid\muh(\val(h))=m\}=-g(m)$ mit
$g(m)=m-s_{\min}(m)$,
$s_{\min}(m)=\min\{s\mid M_s(m{+}1)\ge m!\}$; die
Terminaldarstellung von $x^*=M_{s_{\min}}(m{+}1)$ ist ein Zeuge.
\end{satz}

\begin{proof}
\emph{Schranke.} Sei $h=(n,x)$ mit $\muh(x)=m$, also $x\ge m!$.
Links: $\KV(h)=(m,x)$, und für $P=\operatorname{struct}(x,m)$ gilt
koeffizientenweise $P\preceq M_{\sigma(P)}$, also
$x=P(m{+}1)\le M_{\sigma(P)}(m{+}1)$; wegen $x\ge m!$ folgt
$\sigma(P)\ge s_{\min}$ und damit
$\hor(\KS\KV(h))=\sigma(P)\ge s_{\min}$. Rechts: Für
$P_n=\operatorname{struct}(x,n)$ ist die Auswertung monoton
(nichtnegative Koeffizienten), also
$\val(\KS(h))=P_n(\sigma(P_n){+}1)\le P_n(n{+}1)=x<(m{+}1)!$ und
damit $\hor(\KV\KS(h))\le m$. Zusammen
$\delta(h)\ge s_{\min}-m=-g(m)$.

\emph{Annahme.} Es ist $x^*=M_{s_{\min}}(m{+}1)\in[m!,(m{+}1)!)$
-- untere Grenze nach Definition von $s_{\min}$, obere wegen
$M_{s_{\min}}\preceq M_m$ und $M_m(m{+}1)=(m{+}1)!-1$. Sei
$h^*=(x^*,x^*)$ die Terminaldarstellung. $\KS$ fixiert $h^*$
(konstante Gestalt), also
$\hor(\KV\KS(h^*))=\muh(x^*)=m$. Andererseits ist $M_{s_{\min}}$
im Horizont $m$ eine gültige Ziffernfolge
($c_d\le s_{\min}-d\le m-d$) mit Wert $x^*$, nach der Bijektion
aus Kapitel~\ref{ch:horizonte} also \emph{die} Ziffernfolge:
$\operatorname{struct}(x^*,m)=M_{s_{\min}}$, und
$\hor(\KS\KV(h^*))=\sigma(M_{s_{\min}})=s_{\min}$. Also
$\delta(h^*)=s_{\min}-m=-g(m)$.
\end{proof}

\section*{Zu Kapitel 5: der Seitentreue-Satz}

\begin{satz}[Satz~\ref{satz:seitentreue}, Wert-Gestalt-Fälle]
\emph{(A)} Ist $\oplus$ wertgetragen und $\otimes$ strukturgetragen,
so ist bereits die Links-Distributivität für jede Wahl der
Horizontregeln unerfüllbar. \emph{(B)} Ist $\oplus$ strukturgetragen
und $\otimes$ wertgetragen, so sind Links-Distributivität und
Kommutativität von $\otimes$ zusammen unerfüllbar.
\end{satz}

\begin{proof}
\emph{(A).} Sei $u=(1,1)$ (Struktur $1$), $e_1=u$,
$e_m=e_{m-1}\oplus u$; da $\oplus$ wertgetragen ist, gilt
$\val(e_m)=m$, und die Struktur $Q_m$ von $e_m$ ist irgendeine
gültige Darstellung von $m$. Links-Distributivität liefert induktiv
\begin{equation}\label{eq:app-kette}
\val(a\otimes e_m)=m\cdot\val(a\otimes u).
\end{equation}
\emph{Schritt 1: $Q_m$ ist konstant.} Für $c\ge1$ hat $a_c=(c,c)$
die konstante Struktur $c$ (Terminaldarstellung); Konstanten werten
überall zu sich selbst aus, also $\val(a_c\otimes u)=c$. Mit
\eqref{eq:app-kette} folgt $c\,Q_m(y_c{+}1)=mc$, also
$Q_m(y_c{+}1)=m$, wobei $y_c$, der Horizont von $a_c\otimes e_m$,
wegen $y_c\ge\sigmah(cQ_m)=\max_d(d+c\,q_d)\ge c$ unbeschränkt
wächst ($Q_m\ne0$, denn $\val(e_m)=m\ge1$).
Alle Auswertungspunkte liegen im streng monotonen Bereich
(Argumente $\ge\sigmah+1\ge\deg+2$; fallende Fakultäten sind dort
streng wachsend). Bei $\deg Q_m\ge1$ nähme $Q_m$ den Wert $m$
höchstens einmal an -- Widerspruch. Also ist $Q_m$ konstant, und
mit $a=u$ folgt $Q_m=m$.
\emph{Schritt 2: Konstanten sterben.} Sei $g=(2,3)$ (Struktur $X$),
sei $t_g$ der Horizont von $g\otimes u$ und
$r=\val(g\otimes u)=t_g{+}1\ge\sigmah(X)+1=3$ fest. Nach
Schritt~1 hat $g\otimes e_m$ die Struktur $X\cdot m=mX$ mit
$\sigmah(mX)=m{+}1$; auf seinem Horizont $T_m\ge m{+}1$ ist also
$\val(g\otimes e_m)=m(T_m{+}1)\ge m(m{+}2)$; nach
\eqref{eq:app-kette} ist dieser Wert $mr$ -- Widerspruch für
$m>r-2$.

\emph{(B).} Nun addieren sich Strukturen und multiplizieren sich
Werte. Sei $f_1=u$, $f_m=f_{m-1}\oplus u$; dann hat $f_m$ die
konstante Struktur $m$ und den Wert $m$. Sei $A$ die Struktur von
$g\otimes u$. Induktiv hat $g\otimes f_m$ die Struktur $mA$; als
$\otimes$-Ergebnis vom Wert $3m$, realisiert auf seinem Horizont
$T_m$, gilt $\sigmah(mA)\le T_m$ nach der Gültigkeit der
Darstellung (Ziffernschranke), also $T_m\ge m$, und
$A(T_m{+}1)=3$ an unbeschränkt vielen Stellen. Wie in Schritt~1 ist
$A$ konstant, also $A=3$. Betrachte $f_2'=g\oplus g$ mit Struktur
$2X$, Horizont $H\ge\sigmah(2X)=3$, Wert $2(H{+}1)\ge8$.
Links-Distributivität und Kommutativität geben
$u\otimes(g\oplus g)=(g\otimes u)\oplus(g\otimes u)$ mit Struktur
$2\cdot3=6$ (Konstante), also Wert $6$; links steht aber
$1\cdot2(H{+}1)\ge8$. Widerspruch.
\end{proof}

\begin{satz}[nichtkommutativer Fall]
Ist $\oplus$ strukturgetragen und $\otimes$ wertgetragen und gelten
beide Distributivgesetze, so ist das System unerfüllbar -- ohne
Kommutativität, Assoziativität oder neutrale Elemente.
\end{satz}

\begin{proof}
Die Rechts-Distributivität liefert induktiv
$P_{f_m\otimes a}=m\cdot P_{u\otimes a}$; da $f_m\otimes a$ den
Wert $m\cdot\val(a)$ hat und der Horizont
$\ge\sigma(m\,P_{u\otimes a})\ge m$ wächst, ist $P_{u\otimes a}$
konstant gleich $\val(a)$ (Monotonie-Argument wie oben).
Insbesondere $P_{u\otimes g}=3$; das $g\oplus g$-Finale läuft dann
wörtlich wie im kommutativen Fall über die Links-Distributivität.
\end{proof}

\section*{Zu Kapitel 5 und 8: die dreiseitige Seitentreue}

\begin{satz}[Trilogie]
Bruchgetragene Addition ist nie total; wertgetragene Addition mit
bruchgetragener Multiplikation sowie strukturgetragene Addition mit
bruchgetragener Multiplikation sind für jede Selektorwahl
unerfüllbar (Links-Distributivität genügt).
\end{satz}

\begin{proof}
\emph{Totalität:} $F(u)+F(u)=\tfrac12+\tfrac12=1\notin[0,1)$, und
jedes Element hat Bruch $<1$.

\emph{Wert-$\oplus$, Bruch-$\otimes$:} Ketten $e_m$ mit
$\val(e_m)=m$, Horizont $N_m\ge\mu(m)$, Bruch $m/(N_m{+}1)!$.
Distributivität gibt $\val(u\otimes e_m)=m\,w$ mit
$w=\val(u\otimes u)$ -- ganzzahlig und $\ge1$, denn der Bruch von
$u\otimes u$ ist $\tfrac14>0$. Andererseits ist, mit $h_m$ dem
Horizont von $u\otimes e_m$,
$\val(u\otimes e_m)=\tfrac12\cdot\frac{m}{(N_m+1)!}(h_m{+}1)!$,
also $(h_m{+}1)!/(N_m{+}1)!=2w>1$ konstant -- aber ein
Fakultätsquotient $>1$ ist mindestens $N_m{+}2\to\infty$.

\emph{Struktur-$\oplus$, Bruch-$\otimes$:} Ketten $f_m$ mit
konstanter Struktur $m$, Horizont $N_m\ge m$, Bruch
$m/(N_m{+}1)!$. Distributivität gibt $P_{u\otimes f_m}=m\,A$ mit
$A=P_{u\otimes u}$, $\deg A=:d$; der Bruchvergleich auf dem
Horizont $T_m$ von $u\otimes f_m$ liefert nach Kürzen von $m$
\[
2\,A(T_m{+}1)=(T_m{+}1)^{\underline{\,k_m}},\qquad
k_m=T_m-N_m\ge1 .
\]
Das Produkt rechts ist mindestens $(N_m+k_m/2)^{k_m/2}$, die linke
Seite höchstens $C(N_m{+}k_m{+}1)^d$; wegen $N_m\to\infty$ ist
$k_m\le2d{+}1$ für große $m$, ein Wert $k^*$ tritt unendlich oft
auf, und $2A(Y)=Y^{\underline{k^*}}$ gilt als Polynomidentität.
Dann hätte $A$ den Koeffizienten $\tfrac12$ -- Ziffern sind ganz.
(Allgemeine Multiplikatoren $a$: dieselbe Rechnung mit
$c=(n_a{+}1)!/v_a>1$ statt $2$.)
\end{proof}

\section*{Zu Kapitel 6: Asymptotik und Ziffernformel}

\begin{satz}[Asymptotik der Sprungstellen]
Sei $m_g$ die kleinste Schale mit
$M_{m-g}(m{+}1)\ge m!$, wobei $M_s=\sum_{d<s}(s{-}d)\ff Xd$ das
Maximalpolynom ist. Mit
$c_g=(g{+}2)!\sum_{k\ge1}k/(g{+}k{+}1)!$ gilt
$(g{+}2)!/c_g\le m_g+1\le(g{+}2)!$ und $c_g=1+O(1/g)$, also
$m_g/(g{+}2)!\to1$.
\end{satz}

\begin{proof}
Umindexieren ($k=s{-}d$, $s=m{-}g$) liefert exakt
\[
\frac{M_{m-g}(m{+}1)}{m!}
=(m{+}1)\sum_{k=1}^{m-g}\frac{k}{(g{+}k{+}1)!}=:F(m).
\]
$F$ wächst in $m$ (der Vorfaktor streng, die Summe durch neu
hinzukommende nichtnegative Terme), die Schwelle ist also
wohldefiniert. Für $m{+}1\ge(g{+}2)!$ genügt der Term $k{=}1$:
$F(m)\ge(m{+}1)/(g{+}2)!\ge1$. Umgekehrt ist
$F(m)\le(m{+}1)\,c_g/(g{+}2)!$, also $F(m)<1$ für
$m{+}1<(g{+}2)!/c_g$.
Schließlich ist mit $x=1/(g{+}3)$
\[
c_g\le1+\sum_{k\ge2}k\,x^{k-1}
=1+\frac{x(2-x)}{(1-x)^2}=1+O(1/g). \qedhere
\]
\end{proof}

\begin{satz}[Ziffernformel für die Nicht-Springer]
Sei $d=(d_1,\ldots,d_n)$ die $H_n$-Ziffernfolge von $n!$ und $P$
die erste Position mit maximaler Ziffer ($d_P=P$), bzw.\ $P=n$,
falls keine existiert. Dann gilt
\[
N(n)=\sum_{p=1}^{P}\min(d_p,p)\cdot\frac{n!}{p!}.
\]
\end{satz}

\begin{proof}
$N(n)$ zählt Ziffernfolgen mit $a_i\le i-1$ und Wert $<n!$. Wir
zerlegen nach der ersten Position $p$, an der sich $a$ von $d$
unterscheidet. Der Präfix muss mit $d$ übereinstimmen -- möglich
nur, solange $d_j\le j-1$, also bis zur ersten maximalen Ziffer
einschließlich. An der Stelle $p$ muss $a_p<d_p$ und $a_p\le p-1$
gelten: $\min(d_p,p)$ Möglichkeiten. Der Suffix ist frei innerhalb
der Kappen: $\prod_{j>p}j=n!/p!$ Möglichkeiten, und jeder solche
Suffix bleibt strikt unterhalb, weil der volle (ungekappte)
Suffixraum das Intervall $[0,w_p)$ mit dem Stellengewicht
$w_p=(n{+}1)!/(p{+}1)!$ bijektiv füllt. Jeder Nicht-Springer wird genau
einmal gezählt.
\end{proof}

\begin{satz}[$\beta$-Reduktion]
$\ln\beta(r,r)=\ln\max_{t}\binom{r-1}{t}^2(r{-}1{-}t)!+O(\log r)$.
\end{satz}

\begin{proof}
Untere Schranke: Das Tripel $(p,q,j)=(r{-}1,r{-}1,r{-}1{-}t)$ trägt
zum Grad $r{-}1{+}t$ den Summanden
$\binom{r-1}{t}^2(r{-}1{-}t)!$ bei, und alle Beiträge sind
nichtnegativ. Obere Schranke: Jeder Beitrag ist höchstens
$r^2\binom{r-1}{j}^2j!$, und pro Grad gibt es höchstens $r^3$
Tripel; also $\beta\le2r+r^5\max_t(\cdot)$.
\end{proof}

\begin{satz}[$\beta$-Asymptotik]
$\ln(\beta(r,r)/r!)=2\sqrt r+O(\log r)$.
\end{satz}

\begin{proof}
Nach dem Reduktionssatz genügt die Aussage für den Maximalterm
$\binom st^2(s-t)!/r!=s!/(r\,t!^2(s-t)!)$ mit $s=r-1$. Nach oben:
Stirling-Schranken geben $\ln(s!/(t!^2(s-t)!))\le-h(t)+O(\log s)$
mit $h(t)=2t\ln t-t+(s{-}t)\ln(s{-}t)-s\ln s$; Konkavität liefert
$h(t)\ge t(2\ln t-\ln s-2)$, und diese Funktion hat ihr Minimum
$-2\sqrt s$ bei $t=\sqrt s$. Nach unten: $t=\lceil\sqrt s\rceil$
mit $\binom st\ge(s-t)^t/t!$ ergibt
$t\ln((s-t)/t^2)+2t-O(\log r)=2\sqrt s-O(\log r)$.
\end{proof}

\begin{satz}[Ein-Ziffern-Fakultäten]
$n!$ ist genau dann Einzelziffer in $H_n$, wenn $n=(i{+}1)!/d-1$
mit $1\le d\le i$; die Ziffer ist $d$ an Position $i$, und pro
Position gibt es genau $i$ Fälle.
\end{satz}

\begin{proof}
Einzelziffer $d$ an Position $i$ heißt $n!=d\,(n+1)!/(i+1)!$, also
$d=(i+1)!/(n+1)\le i$. Umgekehrt teilt jedes $d\le i$ die Zahl
$(i+1)!$, und $n+1=(i+1)!/d$ macht die Gleichung exakt; die
Position existiert, denn aus $n{+}1=(i{+}1)!/d\ge(i{+}1)!/i\ge
i{+}1$ folgt $i\le n$, und die eindeutige Ziffernentwicklung
besteht dann aus genau dieser Ziffer.
$d$ durchläuft bijektiv $1,\ldots,i$.
\end{proof}

\section*{Zu Kapitel 6 und 8: Interchange-Positivität, Scharfheit,
Nullzähler}

Für $x<(n{+}1)!$ sei $z_n(x):=P(n{+}2)$ mit
$P=\operatorname{struct}(x,n)$ -- der Wert nach einer
Nullerweiterung. Der \emph{Interchange-Defekt} eines Elements
$h=(n,x)$ ist
$\Omega(h)=z_n(x)-z_{n+1}(x)$: der Wert nach „erst wachsen, dann
werttreu liften“ minus der Wert nach „erst liften, dann wachsen“.

\begin{lemma}[Rekursionsformel]
Für $x=q(n{+}1)+r$, $0\le r\le n$:
$z_n(x)=r+(n{+}2)\,z_{n-1}(q)$.
\end{lemma}

\begin{proof}
Aus der Ziffernrekursion folgt
$P_{n,x}(X)=r+X\cdot P_{n-1,q}(X{-}1)$, denn der Koeffizientenshift
um einen Grad ist wegen
$\ff Xd=X\cdot\ff{(X{-}1)}{d-1}$ die Multiplikation mit $X$ bei
Argumentverschiebung. Auswerten bei $X=n{+}2$.
\end{proof}

\begin{lemma}[strikte Monotonie und Superadditivität]
$z_n(x)<z_n(x{+}1)$, folglich $z_n(y{+}s)\ge z_n(y)+s$.
\end{lemma}

\begin{proof}
Induktion über $n$; Basis $n=1$: $z_1(0)=0<1=z_1(1)$. Ohne
Übertrag wächst $z_n$ um genau 1; beim
Übertrag ($r=n$) ist
$z_n(x{+}1)-z_n(x)=(n{+}2)(z_{n-1}(q{+}1)-z_{n-1}(q))-n
\ge(n{+}2)-n=2$.
\end{proof}

\begin{lemma}[Multiplikationslemma]
Für $c\ge\ell{+}2$ und $cy<(\ell{+}1)!$ gilt
$z_\ell(cy)\ge(c{+}1)\,z_\ell(y)$.
\end{lemma}

\begin{proof}
Induktion über $\ell$; Basisfälle leer bzw.\ trivial. Sei
$y=q(\ell{+}1)+r$ und $cr=s(\ell{+}1)+u$ mit $0\le u\le\ell$; dann
$cy=(cq{+}s)(\ell{+}1)+u$ und nach der Rekursionsformel
$z_\ell(cy)=u+(\ell{+}2)\,z_{\ell-1}(cq{+}s)$. Superadditivität und
Induktionsvoraussetzung (Regime bleibt erhalten:
$c\ge\ell{+}2>(\ell{-}1){+}2$; Gültigkeit: $cq<\ell!$) geben
$z_{\ell-1}(cq{+}s)\ge(c{+}1)z_{\ell-1}(q)+s$. Es bleibt
$u+(\ell{+}2)s\ge(c{+}1)r$: Mit $c=\ell{+}1{+}e$, $e\ge1$, ist
$u=er\bmod(\ell{+}1)$, und die Ungleichung ist äquivalent zu
$r(c{-}\ell{-}1)=er\ge er\bmod(\ell{+}1)=u$ -- wahr.
\end{proof}

\begin{satz}[Interchange-Positivität]
$\Omega(h)\ge0$ für alle $h\in\Hori$.
\end{satz}

\begin{proof}
Nach Definition von $\Omega$ ist zu zeigen:
$z_{n+1}(x)\le z_n(x)$ für $x<(n{+}1)!$. Sei
$x=q'(n{+}2)+r'$. Dann
\[
z_{n+1}(x)=r'+(n{+}3)\,z_n(q')
\le r'+z_n\bigl((n{+}2)q'\bigr)
\le z_n\bigl(q'(n{+}2)+r'\bigr)=z_n(x),
\]
erst mit dem Multiplikationslemma ($c=n{+}2$, exakt an der
Regime-Grenze), dann mit der Superadditivität.
\end{proof}

\begin{satz}[Gleichheitsfälle; Register S32]
$\Omega(h)=0$ gilt genau dann, wenn beide Schritte dieser Kette
scharf sind: das Multiplikationslemma mit $c=n{+}2$ für den
\emph{scharfen Quotienten} $q'$, also
$z_n((n{+}2)q')=(n{+}3)\,z_n(q')$, und die Superadditivität
übertragsfrei, also $z_n((n{+}2)q'+r')=z_n((n{+}2)q')+r'$.
\end{satz}

\begin{proof}
$\Omega(h)$ ist die Summe der beiden Abweichungen
$z_n((n{+}2)q')-(n{+}3)z_n(q')$ und
$z_n((n{+}2)q'+r')-z_n((n{+}2)q')-r'$; beide sind nach
Multiplikationslemma bzw.\ Superadditivität nichtnegativ, ihre
Summe ist also genau dann null, wenn beide null sind.
\end{proof}

\begin{satz}[Ziffernentfaltung der Scharfheit]
Für $c\ge\ell+2$, $cy<(\ell+1)!$, $y=q(\ell+1)+r$, $e=c-\ell-1$
gilt: $z_\ell(cy)=(c{+}1)z_\ell(y)$ genau dann, wenn $er\le\ell$,
$(cq\bmod\ell)+r\le\ell-1$ und dieselbe Bedingung rekursiv für
$(\ell{-}1,c,q)$ gilt (Abbruch: für $\ell=0$ ist die Bedingung
leer).
\end{satz}

\begin{proof}
Die Kette $z_\ell(cy)=u+(\ell+2)z_{\ell-1}(cq+s)$ aus dem Beweis
des Multiplikationslemmas ist genau dann durchgehend scharf: die
arithmetische Abschätzung bei $er\le\ell$ (dann $s=r$, $u=er$),
die Superadditivität bei Übertragsfreiheit der $s=r$
Einzelschritte ab $cq$ -- deren Einerstelle ist $cq\bmod\ell$, die
Bedingung also $(cq\bmod\ell)+r\le\ell-1$ --, die dritte per
Definition rekursiv. Exhaustiv verifiziert für $\ell\le7$.
\end{proof}

\begin{satz}[Nullzähler-Schranke]
$\#\{\Omega=0\text{ auf }H_n\}\le(n{+}1)^{n+1}/n!$, der Anteil
fällt also mindestens wie $e^{-(1+o(1))n\ln n}$.
\end{satz}

\begin{proof}
Nach dem Gleichheitssatz ist jedes solche Element durch $(q',r')$
mit scharfem $q'$ bestimmt, und die Übertragsfreiheit erzwingt
$r'\le n$ (die Einerstelle wächst pro Schritt um eins und muss vor
jedem Schritt unter ihrem Maximum $n$ bleiben). Die Scharfheit
erzwingt nach der Ziffernentfaltung auf Ebene $\ell$ die
Ziffernschranke $r_\ell\le\lfloor\ell/(n{+}1{-}\ell)\rfloor$, und
$\lfloor\ell/j\rfloor+1\le(n{+}1)/j$ mit $j=n{+}1{-}\ell$ liefert
für die $q'$ das Produkt $(n{+}1)^n/n!$; zusammen mit den $n{+}1$
Möglichkeiten für $r'$ folgt die Schranke.
\end{proof}

\section*{Zu Kapitel 9: die Inversionsfolgen-Einordnung}

\begin{satz}[Register S44]
$H_n$ steht in Bijektion zu den Inversionsfolgen der Länge $n{+}1$
und zu $S_{n+1}$. Der Strukturhorizont ist
$\sigma(P_a)=n-\min_{e_j>0}((j{-}1)-e_j)$ für $a\ne0$ (für $a=0$
ist $\sigma=0$) mit $e=(0,a_1,\ldots,a_n)$, und
$|\{a\in H_n:\sigma(P_a)=r\}|=r\cdot r!$ für $1\le r\le n$.
\end{satz}

\begin{proof}
Setzt man $e_1=0$ und $e_{i+1}=a_i$, so ist $0\le e_j<j$ für alle
$j=1,\ldots,n{+}1$ -- die definierende Bedingung einer
Inversionsfolge; die Abbildung $\pi\mapsto(e_j)_j$ mit
$e_j=|\{i<j:\pi_i>\pi_j\}|$ ist die klassische Bijektion zu
$S_{n+1}$, und beide Seiten haben $(n{+}1)!$ Elemente. Für den
Strukturhorizont: $c_d=a_{n-d}=e_{n-d+1}$, also
$d+c_d=(n{+}1-j)+e_j$ mit $j=n-d+1$; Maximieren von $d+c_d$ über
$c_d>0$ ist Minimieren von $(j{-}1)-e_j$ über $e_j>0$, und
$\max(d+c_d)=n-\min_{e_j>0}((j{-}1)-e_j)$. Zählung: $\sigma(P_a)\le r$
bedeutet $a_i\le r-(n-i)$ für alle $i$ mit $a_i>0$, äquivalent
$a\le M$ koordinatenweise für das zu $H_r$ gehörige Maximalprofil;
die Zahl solcher $a$ in $H_n$ ist $(r{+}1)!$ (die Ziffern oberhalb
Position $r$ sind erzwungen null, darunter frei bis zur Schranke).
Also $|\{\sigma\le r\}|=(r{+}1)!$ und
$|\{\sigma=r\}|=(r{+}1)!-r!=r\cdot r!$. (Exhaustiv bestätigt bis
$n=6$, \texttt{experiments\_einordnung.py}.)
\end{proof}
